\par\Needspace{4\baselineskip}
\begin{center}
\hypertarget{native-066-article048}{}\pdfbookmark[2]{48.}{article-048}48.
\end{center}
Damit sich also der Ort eines Himmelskörpers im Raum für jeden beliebigen Zeitpunkt angeben lässt, müssen bei einer elliptischen Bahn die folgenden Größen bekannt sein.

I. Die mittlere Länge für einen beliebig gewählten Zeitpunkt, der \emph{Epoche} heißt. Bisweilen wird auch die Länge selbst so bezeichnet. Meist wählt man als Epoche den Beginn eines Jahres, nämlich den Mittag des 1. Januar in einem Schaltjahr oder den Mittag des vorausgehenden 31. Dezember in einem Gemeinjahr.

II. Die mittlere Bewegung während eines bestimmten Zeitintervalls, etwa während eines mittleren Sonnentages oder während $365$, $365\tfrac{1}{4}$ oder $36525$ Tagen.

III. Die große Halbachse. Sie könnte zwar immer dann weggelassen werden, wenn die Masse des Körpers entweder bekannt ist oder vernachlässigt werden kann, da die mittlere Bewegung die große Halbachse bereits bestimmt (Art. 7). Der Einfachheit halber werden jedoch üblicherweise stets beide angegeben.

IV. Die Exzentrizität.

V. Die Länge des Perihels.

\SourceLeaf{067}VI. Die Länge des aufsteigenden Knotens.

VII. Die Neigung der Bahn.

Diese sieben Größen heißen \emph{Elemente} der Bewegung des Körpers.

Bei der Parabel und der Hyperbel tritt der Zeitpunkt des Durchgangs durch das Perihel an die Stelle des ersten Elements. Für II gibt man die Größen an, die bei diesen Arten von Kegelschnitten der mittleren täglichen Bewegung entsprechen (vgl. Art. 19 und für die hyperbolische Bewegung die Größe $\lambda k b^{-\frac{3}{2}}$ in Art. 23). Bei der Hyperbel können die übrigen Elemente. ebenso beibehalten werden. Bei der Parabel dagegen, deren große Achse unendlich und deren Exzentrizität $=1$ ist, wird anstelle der Elemente III und IV nur der Abstand im Perihel angegeben.

\par\Needspace{4\baselineskip}
\begin{center}
\hypertarget{native-067-article049}{}\pdfbookmark[2]{49.}{article-049}49.
\end{center}
Nach dem üblichen Sprachgebrauch reicht die Bahnneigung, die wir von $0$ bis $180^0$ messen, nur bis $90^0$. Überschreitet der Winkel zwischen der Bahn und dem Bogen $\AscendingNode B$ (Abb. 1) einen rechten Winkel, so gilt der Winkel zwischen der Bahn und dem Bogen $\AscendingNode A$ (der den ersten Winkel zu $180^0$ ergänzt) als Bahnneigung. In diesem Fall muss man hinzufügen, dass die Bewegung \emph{rückläufig} ist (etwa wenn $E\AscendingNode F$ in unserer Abbildung einen Teil der Bahn darstellt), um ihn von dem anderen Fall zu unterscheiden, in dem die Bewegung \emph{rechtläufig} heißt. Die Länge in der Bahn wird dann gewöhnlich so gemessen, dass sie in $\AscendingNode$ mit der Länge dieses Punktes auf der Ekliptik übereinstimmt und in Richtung $\AscendingNode F$ \emph{abnimmt}. Der Ausgangspunkt, von dem aus die Längen entgegen der Bewegungsrichtung gemessen werden, liegt also in Richtung $\AscendingNode F$ ebenso weit von $\AscendingNode$ entfernt wie der Frühlingspunkt von demselben $\AscendingNode$ in Richtung $\AscendingNode A$. Daher ist in diesem Fall die Länge des Perihels gleich der Länge des Knotens, vermindert um den Abstand des Perihels vom Knoten\raisebox{.4ex}{.} Auf diese Weise lässt sich jede der beiden Konventionen leicht in die andere überführen. Wir haben die unsere vorgezogen, damit wir auf die Unterscheidung zwischen rechtläufiger und rückläufiger Bewegung verzichten und für beide stets dieselben Formeln verwenden können, da der übliche Sprachgebrauch sehr häufig zweierlei Rechenregeln erfordert.
